\section{Worked examples}
\label{sec:examples}

The examples are chosen to exercise different layers of the framework.  They
are not substitutes for the full checked Lean developments proposed in
\cref{sec:lean}; they show the complete human proof architecture and the exact
object endpoints.

\subsection{Least divisor as a named bounded certificate}

Let \(z\ge2\).  Define the bounded predicate
\[
  V(z,d)\;:\!\Longleftrightarrow\;
  2\le d\ \wedge\ d\le z\ \wedge\ \exists q\le z\;(d q=z).
\]
All quantifiers are bounded and all atoms are polynomial equalities or order
relations, so \cref{thm:reflection} supplies a Boolean verifier.  The value
\(d=z,q=1\) proves an accepted certificate below \(z+1\).  Apply
\cref{thm:least-code} to obtain the least accepted \(d^*(z)\).

The object theorem first proves that \(d^*\) divides \(z\) and is at least two.
To show primality, suppose \(d^*=ab\) with \(1<a<d^*\) and \(1<b\).  Then
\(a\) divides \(d^*\) and hence divides \(z\); the quotient can be composed
from the witnesses.  Thus \(a\) is an accepted divisor smaller than \(d^*\),
contradicting leastness.  Therefore \(d^*\) is prime.

The public term is justified because the certificate space has the explicit
bound \(z+1\).  If the theorem needed only ``some prime divides \(z\),'' the
certificate could instead be sealed after the primality proof.

\subsection{Rational bisection for the inverse of \texorpdfstring{\(y^3+y\)}{y cubed plus y}}

Fix an integer or rational input \(x\) and define
\(f_x(y)=y^3+y-x\).  The function is strictly increasing because
\[
 f_x(v)-f_x(u)
  =(v-u)(v^2+uv+u^2+1),
\]
and the second factor is positive when \(u<v\).  Choose rational endpoints
\(L_0,U_0\) with \(f_x(L_0)\le0\le f_x(U_0)\).  Such crude bounds are easy to
compute from \(|x|+1\).

At stage \(t\), let \(M_t=(L_t+U_t)/2\).  Exact rational sign comparison is a
bounded integer calculation after clearing positive denominators.  Set
\[
 (L_{t+1},U_{t+1})=
 \begin{cases}
   (M_t,U_t),& f_x(M_t)\le0,\\
   (L_t,M_t),& f_x(M_t)>0.
 \end{cases}
\]
Encode the pair of endpoints and common denominator in one state term.  The
reflected branch bit makes the step primitive recursive.

Use the quantifier-free invariant
\[
  f_x(L_t)\le0\ \wedge\ 0\le f_x(U_t)
   \ \wedge\
  2^t(U_t-L_t)=U_0-L_0.
\]
The base case is the initial bracket.  In the successor case, monotonicity and
the reflected sign branch preserve the bracket, while elementary rational
arithmetic proves the width equation.  \cref{thm:trace} gives the invariant at
an arbitrary query \(k\).  The midpoint or interval is therefore an
approximation to the unique inverse value, with an explicit error
\((U_0-L_0)2^{-k-1}\).

This example has no noncanonical certificate: the trace is deterministic.  It
exercises reflected branching, a packed state, query-first execution, and a
quantifier-free invariant.  A regular name follows by choosing a computable
number of bisection steps for each requested precision.

\subsection{Root clusters instead of unstable branch labels}

Suppose a finite algebra certificate produces disks
\(D_1,\ldots,D_m\) with multiplicities whose union contains all roots of a
polynomial and whose boundary tests certify the root counts.  At a collision,
individual labels are unstable, but the following observations are stable:

\begin{itemize}
\item the multiset of cluster centers and radii modulo permutation;
\item the multiplicity of each connected cluster;
\item symmetric polynomials of roots in a cluster;
\item the product factor associated with a cluster.
\end{itemize}

Let \(V(p,c)\) verify a cluster certificate and \(O(p,c)\) be one of these
observations.  A refinement theorem proves that any two valid certificates
have a common refinement and that their observations agree modulo permutation
or lie within a stated error.  \cref{thm:observed-relation} then exports a
coherent relational observation without choosing a branch.

If a separation residual
\[
  \operatorname{dist}(D_i,D_j)\ge\delta>0\quad(i\neq j)
\]
is supplied, request enough coefficient precision to verify the gaps and order
the disks lexicographically.  A bounded selector may then name each branch.
The label theorem is explicitly conditional on \(\delta\); it does not pretend
to pass continuously through a collision.

\subsection{Finite-net approximate minimization}

Let \(K\) be represented at precision \(r\) by a finite list
\(N_r=(x_0,\ldots,x_{M-1})\), a coverage relation saying every point of \(K\)
is within \(2^{-r}\) of some \(x_i\), and let \(f\) have a supplied modulus
\(\omega\).  Evaluate rational approximations \(v_i\) with error at most
\(2^{-s}\), and select the least index \(i^*\) minimizing the finite value
codes.

The finite search and comparison proof gives
\(v_{i^*}\le v_i\) for all net indices.  For any \(x\in K\), open the coverage
certificate to obtain a nearby \(x_j\).  The modulus gives
\(|f(x)-f(x_j)|\le\omega(2^{-r})\), and approximation errors compare
\(f(x_j)\) with \(v_j\) and \(v_{i^*}\) with \(f(x_{i^*})\).  Hence
\[
  f(x_{i^*})\le f(x)+2^{1-s}+\omega(2^{-r}).
\]
Seal the noncanonical nearby-net witness.  The public endpoint names the least
finite minimizer only because the list is explicitly bounded.  The coverage
and modulus formulas remain residuals unless separately produced.

\subsection{Conditional complete enumeration of Diophantine points}

Let \(P(\vec x)=0\) be an integer polynomial equation and let \(h(\vec x)\) be
a coded height.  Suppose an external mathematical argument supplies the object
producer
\[
  \Gamma\turnstilep
  \forall\vec x\,[P(\vec x)=0\to h(\vec x)\le H].
  \tag{HB}
\]
There is a primitive-recursive enumeration of every normalized integer or
rational tuple of height at most \(H\).  Filter it using reflected polynomial
equality and remove duplicates by reflected code equality.  Let \(L_H\) be the
resulting list.

A finite trace proves that each retained element is a solution and each
rejected enumerated code is either malformed or not a solution.  Given an
arbitrary solution \(x\), apply (HB) to place it below the enumeration bound;
normalization yields its unique code, and the filter proof shows that code is
retained.  Thus
\[
  \mathsf{HeightBound}(P,H),\Gamma
    \turnstilep \mathsf{CompleteList}(P,L_H).
\]

The finite enumeration theorem is unconditional; completeness is conditional
on the exact height producer.  Mere finiteness of the solution set cannot
replace (HB).

\subsection{A whole-path example}

Consider three finite stages.  At stage zero, choices \(a\in\{0,1\}\) are both
locally valid.  From \(a=0\), the only stage-one choice eventually dead-ends;
from \(a=1\), an accepted terminal path exists.  A greedy ``least valid
choice'' incorrectly selects \(0\).

Backward viability marks the terminal accepted nodes, then marks exactly the
stage-one nodes with accepted successors, and finally marks only \(a=1\) as
viable.  Forward least selection among \emph{viable} successors constructs the
accepted path.  The object proof is a fixed finite instance of the viability
theorem.  This tiny example captures the general reason separate existential
witnesses cannot simply be concatenated.
